\documentclass{article}
\usepackage[utf8]{inputenc}
\usepackage[T1]{fontenc}
\usepackage{geometry}
\usepackage{graphicx}
\usepackage{float}
\usepackage{hyperref}
\usepackage[english,indonesian]{babel}

\title{Optimasi Pemrosesan Template LaTeX Menggunakan Pendekatan Berbasis AI}
\author{Academic Support Agent}
\date{\today}

\begin{document}
\maketitle

\begin{abstract}

\end{abstract}

\chapter{PENDAHULUAN}


\chapter{TINJAUAN PUSTAKA}


% Original content (commented out):
% %---------------------------------------------------------------------------------------|
%	CONFIG
%---------------------------------------------------------------------------------------|
\input{structure.tex}
\input{cover.tex}
%---------------------------------------------------------------------------------------|


%---------------------------------------------------------------------------------------|
%	TITLE
%---------------------------------------------------------------------------------------|
\title{Lorem ipsum dolor sit amet, consectetur adipiscing elit. Ut semper mi felis, et vestibulum risus}
\professor{XXXXXXXXXX XXXXXXXX}
\signature{XXXXX XXXXXXXXXX}
\scode{IS-XXX}
\author{XXXXX XXXXX, XXXXX XXXXX.  \par XXXXX XXXXX, XXXXX XXXXX.  \par XXXXX XXXXX, XXXXX XXXXX.}
\acode{27XXXXXX \par 27XXXXXX \par 27XXXXXX}
\date{2019}
\begin{document}
    \maketitle
    \tableofcontents
    \newpage
%---------------------------------------------------------------------------------------|


%---------------------------------------------------------------------------------------|
%	DOCUMENT
%---------------------------------------------------------------------------------------|
    \section{Introduction}

%---------------------------------------------------------------------------------------|
\end{document}
%---------------------------------------------------------------------------------------|

\end{document}